% NFC chapter. Section order is the per-peripheral template enforced by
% python/check_intro_names.py. Vector numbers are frozen by the generator
% (python/generate.py, _mcuMpIrqVectors).
\section{Overview}

The NFC controller (\peripheral{NFCx}) is an ISO/IEC 14443 Type A tag emulation engine. It presents a passive contactless tag to an external reader: it decodes the reader-to-tag frames, runs the tag transaction state machine, and load-modulates the tag's response. The 13.56 MHz analog front end is off-die; the block is digital only and connects to an external front end through a small demodulated interface.

\begin{itemize}
	\item ISO 14443A tag emulation: REQA/WUPA to ATQA, bit-frame anticollision with a 4-byte UID to SAK, then Type-2 READ answered from a firmware-filled payload window (\bitfield{NFCAUTOREAD})
	\item Miller decoding of reader frames; Manchester subcarrier (fc/16) load modulation of tag responses
	\item Byte framing with odd parity and CRC\_A (reflected \texttt{0x8408}, initial value \texttt{0x6363})
	\item Provisioned 4-byte UID (\register{NFCxUID}), programmable ATQA and SAK (\register{NFCxCFG}), hardware BCC
	\item Two 64-byte indexed windows behind one index/data pair (\register{NFCxIDX}, \register{NFCxDATA}): the payload window and the received-frame buffer
	\item Programmable protocol timing (\register{NFCxTIM}): bit period, subcarrier divisor, frame delay time
	\item Registered read path: no register read has a side effect
	\item Four interrupt sources, vectors 94 to 97
\end{itemize}

\section{Functional description}

\subsection{Analog front-end boundary}

The block never touches the carrier. Its interface to the external front end is a demodulated receive envelope (1 = field present, 0 = a reader pause), an asynchronous field-present level, the carrier-derived clock \texttt{rf\_clk} that clocks the protocol core, a load-modulation drive output, a transmit-active window, and a listen-power enable. These wires are alternate functions on port 6 (Section \ref{s:pinsConfig}).

The block spans three clock domains: the bus clock for the register interface, a free-running SMCLK reference that hosts the synchronizers and the write-1-to-clear retirement, and \texttt{rf\_clk}. Quasi-static configuration (UID, ATQA, SAK, timing) is latched per transaction so that a multi-byte value is never sampled while software is updating it, and received-frame bytes cross by a flag-gated handshake: the buffer is written before \bitfield{NFCRXFRAMEF} is raised, so a handler that sees the flag can read the frame.

\subsection{Transaction engine}

With \bitfield{NFCEN} set the protocol core runs; with \bitfield{NFCLISTEN} also set it answers readers. The engine walks the 14443A sequence on its own: REQA or WUPA is answered with the ATQA from \register{NFCxCFG}, the anticollision cascade with the UID from \register{NFCxUID} and its hardware-derived BCC, and SELECT with the SAK. With \bitfield{NFCAUTOREAD} set (its reset value), a Type-2 READ is answered directly from the payload window without firmware. Any other command, or every command when \bitfield{NFCAUTOREAD} is clear, is surfaced to firmware: \bitfield{NFCRXFRAMEF} is set, \register{NFCxRXST} reports the command byte, length and CRC/parity result, and the raw frame is readable through the frame buffer. Firmware answers with \register{NFCxTXCTL}: it fills the payload window, writes the length to \bitfield{NFCTXLEN}, selects CRC appending with \bitfield{NFCTXAPPCRC}, and sets \bitfield{NFCTXGO}.

\subsection{Indexed windows}

Both 64-byte windows are reached through one index/data pair: \bitfield{NFCIDXSEL} in \register{NFCxIDX} selects the payload window (0) or the received-frame buffer (1), \bitfield{NFCIDX} selects the byte, and each access to \register{NFCxDATA} reads or writes it. With \bitfield{NFCIDXAINC} set the index advances after every access, so a record is streamed with consecutive data accesses.

\subsection{Protocol timing}

\register{NFCxTIM} carries the bit period (\bitfield{NFCETU}), the subcarrier half period (\bitfield{NFCSUBCDIV}) and the frame delay time (\bitfield{NFCFDT}), all counted in \texttt{rf\_clk} ticks. The reset values follow the 13.56 MHz grid; a simulation bench writes them small to compress protocol time.

\section{Interrupts and events}

The four flags in \register{NFCxSR} are write-1-to-clear and each has an enable in \register{NFCxCR}. Table \ref{t:nfc-irq} lists them. The field-detect level is also available to the boot gate in \peripheral{PWRCTRL} as a wake source, which does not use a vector.

\begin{table}[htb]
	\centering
	\caption{\peripheral{NFC0} interrupt vectors.}
	\label{t:nfc-irq}
	\begin{tabular}{ l p{4.6cm} l l }
	\textbf{Vector} & \textbf{Condition} & \textbf{Set by} & \textbf{Cleared by} \\ \hline
	94 & RF field detected & \bitfield{NFCFIELDF} & Writing 1 to \bitfield{NFCFIELDF} \\
	\rowcolor{tablehighlightcolor}
	95 & Reader frame received for firmware & \bitfield{NFCRXFRAMEF} & Writing 1 to \bitfield{NFCRXFRAMEF} \\
	96 & Tag response transmit done & \bitfield{NFCTXDONEF} & Writing 1 to \bitfield{NFCTXDONEF} \\
	\rowcolor{tablehighlightcolor}
	97 & Receive CRC or parity error & \bitfield{NFCCRCERRF}, \bitfield{NFCPARERRF} & Writing 1 to the flag \\
	\hline
	\end{tabular}
\end{table}

\section{Programming}

\subsection{Provisioning and arming the tag}

\begin{enumerate}
	\item Write \register{SYSCLKCR} so that SMCLK runs from HFXT (Section \ref{ss:clocks}).
	\item Set the \register{PxSEL} bits and \register{PxAFS} fields of the six front-end pins on port 6.
	\item Write the UID to \register{NFCxUID}.
	\item Write \register{NFCxCFG}.\bitfield{NFCATQA} and \register{NFCxCFG}.\bitfield{NFCSAK}.
	\item Write \register{NFCxIDX} with \bitfield{NFCIDXSEL} = 0, \bitfield{NFCIDX} = 0 and \bitfield{NFCIDXAINC} = 1.
	\item Write each payload byte to \register{NFCxDATA}.
	\item Set the interrupt enables required in \register{NFCxCR}.
	\item Set \register{NFCxCR}.\bitfield{NFCEN} and \register{NFCxCR}.\bitfield{NFCLISTEN}.
\end{enumerate}

\subsection{Servicing a frame delivered to firmware}

\begin{enumerate}
	\item Read \register{NFCxRXST} for the command, length and CRC/parity result.
	\item Write \register{NFCxIDX} with \bitfield{NFCIDXSEL} = 1 and \bitfield{NFCIDXAINC} = 1, then read the frame bytes from \register{NFCxDATA}.
	\item Fill the payload window with the response.
	\item Write \register{NFCxTXCTL} with \bitfield{NFCTXLEN}, \bitfield{NFCTXAPPCRC} and \bitfield{NFCTXGO}.
	\item Write 1 to \register{NFCxSR}.\bitfield{NFCRXFRAMEF}.
\end{enumerate}

\section{Usage notes}

\paragraph{Keep SMCLK on HFXT while the tag is armed.} The synchronizers and flag retirement run on SMCLK; a stopped or slow SMCLK stalls them.

\paragraph{Use a byte store to arm LISTEN when the payload is auto-read.} A full-word write of \register{NFCxCR} also writes the byte lane holding \bitfield{NFCAUTOREAD}; writing 0 there disables the hardware answer to READ.

\paragraph{Provision before enabling.} The UID, ATQA, SAK and timing are latched per transaction, so a value written while a reader is mid-transaction takes effect at the next one.

\paragraph{Do not read the block through an alias when it is absent.} In a configuration without the NFC block its address range aliases the \peripheral{MUTEX} page, and a read there claims a mutex.
